%% kapa: exam=1 round=19 year=2008 period=2 form=A subject=economics subject_ko=경제원론 seq=19 number=19
%% source: Q-Net 감정평가사 기출문제 (cst00309) artlSeq=1001962 file=extracted/1008/1001962/2008년 제19회 감정평가사 1차시험 기출문제/2교시.hwp
%% conversion: kordoc 4.24.51 -> tools/gen1.py (auto); review=proofread-sonnet; flags=none
\begin{question}{19}
고정환율제도와 이 제도에서 나타날 수 있는 현상에 대한 설명으로 옳은 것을 모두 고르면?
\begin{framedtext}
ㄱ. 국제수지\textsuperscript{*} 흑자가 발생할 경우 국내 통화공급이 감소한다. \par
ㄴ. 국제수지\textsuperscript{*} 적자가 발생할 경우 중앙은행이 외환을 매각해야 한다. \par
ㄷ. 고정환율제도는 해외에서 발생한 충격을 완화시켜 주는 역할을 한다. \par
ㄹ. 국내 정책목표를 달성하기 위한 통화정책이 제약을 받는다. \par
\end{framedtext}
\textsuperscript{*}국제수지 = 경상수지＋자본수지

\begin{choices}
  \item ㄱ, ㄷ
  \item ㄱ, ㄹ
  \item ㄴ, ㄷ
  \item ㄴ, ㄹ
  \item ㄱ, ㄴ, ㄹ
\end{choices}
\end{question}
\answer{4}
